\section{Bridges to other areas, and the directions they open}\label{sec:overlaps}

The research behind this paper explored many directions quickly, and in the author's view the bridges it built are among its most valuable outcomes. This section describes them. It does not claim that any of them is a complete research programme, and it does not treat a direction that is still incomplete as invalid: much of mathematics consists of programmes that are incomplete and are nevertheless valid mathematics. The audit did not have the compute to examine every direction to the depth needed to decide how promising it is. For each direction the section states what is established, with references to the proved results; what is not verified here; and my assessment.

The assessment addresses two separate questions. The first is the direction's interest on its own merit, for example as a bridge between two areas. The second is its bearing on the Erdős--Straus problem. When a direction relates the equation to a structure that does not take primes as input, the second question is what typed map connects the two, and how far that map is established. The assessments are mine and are marked as such; they are views held at this stage and may change as the directions are explored further.

\subsection{The split-zero support language}\label{ssec:splitzero}

The notes describe shells, classes and certificates through the split-zero semiring $G(R)=R\sqcup\{\tau\}$. In it, $\tau$ is an additive identity distinct from the ring zero $0_R$, and $\tau$ absorbs multiplication. The same semiring is the coefficient algebra of the zeta programme \cite{ZetaReader}, and it is archive §21. The note \emph{A split-zero repair dossier} develops the algebra of $G(R)$: ideals, prime ideals, spectrum, congruences, localisation and semimodules. Its classification statements were checked in a reading pass. One theorem of the dossier needs an additional relation, and the correction keeps to the construction's own convention of keeping $\tau$ and $0_R$ apart.

\begin{proposition}[the dossier's Theorem~7.2, with its relation]\label{prop:dossier}
Let $S=G(R)$, let $M_S$ be its multiplicative monoid, and let $\mathcal R$ be the congruence on the free commutative semiring $\NN[M_S]$ generated by $[a]+[b]\sim[a+b]$ for $a,b\in R$ and $[x]+[\tau]\sim[x]$ for $x\in S$.
\begin{enumerate}[label=(\alph*)]
\item The class of the empty sum $0$ in $\NN[M_S]/\mathcal R$ is $\{0\}$.
\item The natural map $\bar\Theta:\NN[M_S]/\mathcal R\to S$ sends both $0$ and $[\tau]$ to $\tau$, so it is not injective.
\item After adjoining the relation $[\tau]\sim0$, $\bar\Theta$ is an isomorphism.
\end{enumerate}
\end{proposition}
\status{Proved here; found by a reading pass of version~3, new as a proposition in version~4.}
\begin{proof}
(a) Every generating relation relates two non-empty formal sums. Adding a sum, or multiplying by a non-zero sum, preserves non-emptiness, and multiplying by $0$ gives $0\sim0$. So no chain of relations connects $0$ to a non-empty sum.
(b) A semiring homomorphism sends $0$ to $0_S=\tau$, and $\bar\Theta([\tau])=\tau$.
(c) With $[\tau]\sim0$, every formal sum reduces to $[\tau]$ or to $[r]$ for some $r\in R$: the terms $[\tau]$ can be removed by $[x]+[\tau]\sim[x]$, the empty sum is $\sim[\tau]$, and the remaining terms $[r_1],\dots,[r_k]$ combine to $[r_1+\dots+r_k]$. These representatives are pairwise inequivalent, because $\bar\Theta$ sends them to the distinct elements $\tau$ and $r\in R$ of $S$ (recall $\tau\ne0_R$). So $\bar\Theta$ is bijective, and it is a homomorphism by construction.
\end{proof}

The support states also have a typed meaning in the shell arithmetic. For a shell $a$, the three states $\tau$, $0$ and $+$ correspond to three situations:
\begin{itemize}
\item neither $-1$ nor $-4$ lies in the group $H_a$, and the barriers of Section~\ref{ssec:jacobi} apply;
\item $-1$ or $-4$ lies in $H_a$, but no divisor of $a^2$ reaches a target;
\item the shell is occupied.
\end{itemize}
By Proposition~\ref{prop:groupconverse}, under the exponent condition the middle state does not occur. Applied to classes of primes, as in Remark~\ref{rem:termcore}, the same vocabulary calls a class supported when some certificate applies on a lift of it.

My assessment: I think this vocabulary is useful for the problem, because it separates the two reasons for which a shell can be empty. One is group-theoretic, and the barriers handle it; the other is a matter of size, in which the exponents of $a$ are too small to reach a target inside the group. They call for different tools. The algebra of $G(R)$ developed in the dossier is also an object of independent interest.

\subsection{Mock and quantum modularity}

What is established is in Section~\ref{ssec:mock}: the notes' explicit mixed mock functions are Ramanujan's order-$7$ mock theta functions, and the first is $q^{1/2}\widehat Z_0(-\Sigma(2,3,7))$; their shadow carries the $M(2,7)$ representation after conjugation and the $\eta^3$ twist (Theorem~\ref{thm:m27}); and the order-$5$ block carries the Fibonacci data (Proposition~\ref{prop:fibonacci}). Not verified here: the note's support isomorphism from the cubic middle of the shell problem to the three support labels, its mass-gap coordinate for the cubic middle, and its Virasoro support-gap corollary.

My assessment: on its own merit, I think this is the most developed of the structural directions. It has produced a result, Theorem~\ref{thm:m27}, that stands independently of the Erdős--Straus problem and connects Ramanujan's order-$7$ functions, the $3$-manifold invariant of $-\Sigma(2,3,7)$ and the $(2,7)$ minimal model. For its bearing on the equation, the typed map needed runs from the shell data, which depend on $p$ through the factorisation of $(p+R)/4$, to the modular object. The notes' map goes through the finite support labels. An analytic object in the notes whose coefficients depend on that factorisation is the pair-count series of Proposition~\ref{prop:pairzeta}. A typed map between the pair-count series and the mock vector would carry information in both directions, and constructing one is Question~\ref{q:modularobject}.

\subsection{Moonshine, Niemeier lattices and Ogg's primes}

Checked in reading passes (Section~\ref{ssec:structfacts}): the Niemeier lattice with root system $A_5^4D_4$, its umbral group $GL_2(3)$, the $\Gamma_0(6)$ Hauptmodul and the series of the class $6B$. A constructive-algebra preprint of the notes also connects the extension prime $29$ with the Heegner discriminant $163$. The factorisation it uses is classical and correct: $j\bigl((1+\sqrt{-163})/2\bigr)=-640320^3$ with $640320=2^6\cdot3\cdot5\cdot23\cdot29$; by contrast, $5280=2^5\cdot3\cdot5\cdot11$ for the discriminant $67$. Not verified here: the preprint's rank-$15$ Ogg module, its genus-zero and supersingular realisations, and its phase decomposition, apart from the map treated next.

\begin{proposition}[the torsor map]\label{prop:torsor}
The map $\Theta:(\ZZ/18)^\times\to S_{840}$ of \emph{Terminal residual coordinates} §10, given by $1\mapsto1$, $5\mapsto289$, $13\mapsto361$, $17\mapsto169$, $7\mapsto121$, $11\mapsto529$, is a bijection but not a homomorphism. The map $5^k\mapsto289^k$ is an isomorphism, and it differs from $\Theta$ exactly by exchanging the images of $7$ and $13$.
\end{proposition}
\status{Proved here; new as a proposition in version~4.}
\begin{proof}
Both groups are cyclic of order $6$: $5$ generates $(\ZZ/18)^\times$, with powers $5,7,17,13,11,1$, and $289$ generates $S_{840}$, with powers $289,361,169,121,529,1$ modulo $840$. So $5^k\mapsto289^k$ sends $5,7,17,13,11,1$ to $289,361,169,121,529,1$. It agrees with $\Theta$ except at $7$ and $13$, whose images are exchanged. In particular $\Theta(5)^2=289^2\equiv361\pmod{840}$, while $5^2\equiv7\pmod{18}$ and $\Theta(7)=121$.
\end{proof}

The note's odd Ogg phase decomposition (its Theorem~10.1) uses only that $\Theta$ is a bijection, so it holds as stated; with the isomorphism in place of $\Theta$, two of its sectors exchange their contents (a reading pass).

My assessment: the moonshine data organise the residue layer of the problem, meaning the six hard classes and the core, in a way that matches known structures: the umbral group of $A_5^4D_4$, the $\Gamma_0(6)$ Hauptmodul and Ogg's primes. I think this match is of independent interest as an observation. Its bearing on the hard classes depends on the typed map that is used. What is proved is that polynomial identities and fixed-divisor certificates do not reach the square classes (Section~\ref{sec:problem}, Proposition~\ref{prop:nosquare}). I have not seen a map from the shell data to the moonshine data in the notes, and I do not have a view yet on whether one exists.

\subsection{Lorentzian and Apollonian geometry}\label{ssec:lorentz}

The Descartes form $F(x)=2\sum x_i^2-(\sum x_i)^2$ has signature $(3,1)$. The Apollonian group is a subgroup of $O_F(\ZZ)$ of infinite index that is Zariski dense in $O_F$, a thin group, and it acts on hyperbolic $3$-space \cite{SarnakLetter,GLMWY}. The notes identify a Lorentz form of their shell geometry with the Descartes form. The same Lorentzian geometry of circle configurations underlies Kocik's circle-geometric reading of the Koide relation \cite{Kocik}.

The note \emph{Lorentzian fixed fibers, Descartes sheets, and origin-resolved positivity in residual Erdős--Straus shells}, in its section \emph{Finite-dimensional modular flow as a Lorentz boost}, states that the congruence action $X\mapsto A_\rho XA_\rho$, with $A_\rho=\operatorname{diag}(e^{\rho/2},e^{-\rho/2})$, is a Lorentz boost on the Hermitian $2\times2$ matrices with the form $\det$. This is correct. The relation between this boost and the modular flow of the section's title is the following.

\begin{proposition}[modular flow and boost]\label{prop:flowboost}
Let $H_r=\operatorname{diag}(e^r,e^{-r})$ and $X=\begin{psmallmatrix}h+w&u+iv\\u-iv&h-w\end{psmallmatrix}$, with the Minkowski form $\det X=h^2-u^2-v^2-w^2$. For $z\in\C$ the map $X\mapsto H_r^zXH_r^{\bar z}$ lies in $SO^+(1,3)$.
\begin{enumerate}[label=(\alph*)]
\item At $z=it$ it is the modular automorphism $\sigma_t=\operatorname{Ad}(H_r^{it})$ of the state $\varphi_{H_r}$ with density proportional to $H_r$. It acts as the rotation by the angle $2rt$ in the $(u,v)$-plane and fixes $h$ and $w$.
\item At real $z=s$ it is the boost of rapidity $2rs$ in the $(h,w)$-plane and fixes $u$ and $v$. The note's $A_\rho$ is the case $s=\rho/(2r)$.
\end{enumerate}
So the modular flow and the note's boost are the compact and the non-compact real forms of one complex one-parameter subgroup $z\mapsto H_r^z$ of $SL_2(\C)$, acting on Minkowski space by $X\mapsto gXg^*$.
\end{proposition}
\status{Proved here; new in version~4. The coordinate formulas are checked symbolically in \note{misc\_checks.py}.}
\begin{proof}
$\det H_r^z=1$ and $(H_r^z)^*=H_r^{\bar z}$, so the map is the standard action of $SL_2(\C)$ on Hermitian matrices, which lies in $SO^+(1,3)$ because $SL_2(\C)$ is connected. At $z=it$, $H_r^{\bar z}=H_r^{-it}$, so the map is $\operatorname{Ad}(H_r^{it})$, the modular group of $\varphi_{H_r}$; it multiplies the off-diagonal entry $u+iv$ by $e^{2irt}$ and fixes the diagonal, which gives $(h,\,u\cos2rt-v\sin2rt,\,u\sin2rt+v\cos2rt,\,w)$. At $z=s$ it multiplies the diagonal entries $h+w$ and $h-w$ by $e^{2rs}$ and $e^{-2rs}$ and fixes $u+iv$, which gives $(h\cosh2rs+w\sinh2rs,\,u,\,v,\,w\cosh2rs+h\sinh2rs)$.
\end{proof}

My assessment: on its own merit, I think the Descartes setting is a natural place to ask how such Lorentz one-parameter families interact with the Apollonian thin group, and the question is one about thin groups in its own right. A companion note answers it \cite{ApollonianNote}. The orientation-preserving subgroup $A^+$ of the Apollonian group meets every compact subgroup of $SO_F(\R)$ only in the identity, so under any identification of the Descartes space with the Hermitian matrices no nontrivial element of the Apollonian group is a modular automorphism of a faithful state (the modular automorphisms lie in compact subgroups of $SO^+(1,3)$, and the reflections $S_i$ are not in $SO_F$). Boosts occur in the stabilisers of the circles of a packing, which are congruence groups (in the strip packing, the image of $\Gamma^0(4)$ in $PSL_2(\ZZ)$), and also as boosts that fix no circle, such as $S_1S_2S_1S_3S_4S_3$; every other loxodromic element of $A^+$ is a screw motion whose rotation angle is an irrational multiple of $2\pi$. For the bearing on the equation, the notes' map from shell data to Descartes quadruples is not verified here. I have not identified a typed map from shell occupancy to the Apollonian group, and I have no view yet on whether one exists.

\subsection{Additive combinatorics: Kneser's theorem}

The notes apply Kneser's theorem on sumsets in abelian groups \cite{Kneser} to the residues of divisors, producing Star--Kneser forcing, which the archive states as a per-shell sufficient condition for occupancy (Theorems~15.6 and~20.7; located, not verified here). Its converse fails at $p_*$, $R=107$, as the archive records (Proposition~15.7).

My assessment: I think additive combinatorics is a natural toolset for the second reason for which a shell can be empty (Section~\ref{ssec:splitzero}), where the target lies in $H_a$ but may not be reached. In additive notation the residues of the divisors of $a^2$ form a sumset of arithmetic progressions, one for each prime $\ell\mid a$ with $2v_\ell(a)+1$ terms, so sumset growth is exactly what decides whether they fill $H_a$.

\subsection{Computability: the Busy-Beaver notes}

The statement that $4/n=1/x+1/y+1/z$ is solvable for every $n\ge2$ is $\Pi_1$. Each instance is decided by a bounded search: in a solution with $x\le y\le z$ one has $x\le3n/4$, then $y\le2/(4/n-1/x)$, and $z$ is determined. So there is a Turing machine that halts if and only if the conjecture fails, and if it has $k$ states, the value of the Busy-Beaver function at $k$ decides the conjecture. This is the standard bridge between $\Pi_1$ statements and the Busy-Beaver function; such machines have been constructed explicitly for Goldbach's conjecture and for the Riemann hypothesis \cite{Aaronson}. The archive uses Busy-Beaver values, quoted from the literature and not checked here ($S(4)=107$, $\mathrm{Space}(5)=12289$), to select the shells of its §§3 and 20, and its §25 maps Busy-Beaver extrema to shell carriers (located). The proposed pattern in residues modulo $107$ is discussed in Section~\ref{ssec:structfacts}.

The $\Pi_1$ bridge is exact, and the least number of states of a machine that halts exactly when the conjecture fails is a well-posed target (Question~\ref{q:pi1}).

\subsection{Further directions in the same collections}

The collections also contain notes on Cayley--Dickson algebras and sedenion zero divisors, on the Koide relation, on a mass gap and on anomalies, and a split-zero zeta note. The Cayley--Dickson statements were checked in a reading pass (Section~\ref{ssec:structfacts}). The Koide, mass-gap and anomaly notes lie outside the scope of this paper; they were not read for it and are not verified here.

\subsection{The author's other workbenches}\label{ssec:wbedges}

Material crosses between this workbench and the author's other workbenches. The audit note \note{47\_} records every crossing that the audit found as a typed edge, of one of four types:
\begin{itemize}
\item \emph{P}: a proof-level import;
\item \emph{D}: a dependency or citation;
\item \emph{I}: the same identity or object on both sides;
\item \emph{A}: an analogy, with no typed map yet.
\end{itemize}
The table lists the edges that touch this workbench, and adds three that this paper found in the archive (marked \emph{new}).

\begin{center}\small
\begin{longtable}{@{}>{\raggedright\arraybackslash}p{0.13\linewidth}>{\raggedright\arraybackslash}p{0.44\linewidth}>{\raggedright\arraybackslash}p{0.07\linewidth}>{\raggedright\arraybackslash}p{0.28\linewidth}@{}}
\toprule
Edge & What crosses & Type & Status\\
\midrule
\endhead
NS $\to$ ES & The auxiliary torus endomorphism $J=\begin{pmatrix}3&1\\1&5\end{pmatrix}$ (determinant $14$) of the Navier--Stokes manuscript. It enters 21 statements of the supplement (§8) as an operator only. The core identity is $\{(u,v)\in G^2:u^3v=1=uv^5\}=\{(u,u^{-3}):u^{14}=1\}$ & P & audited (\note{21\_} §1); no fluid theorem is used\\
S$^6\to$ ES & The integral monodromies $A_1,A_2,A_\infty$ of the $(3,4,\infty)$ period system, in ES-09 & P & audited for odd levels; the orbits at every level are Theorem~\ref{thm:orbits} (§\ref{ssec:es-s6})\\
ES $\leftrightarrow$ S$^6$ & The Niemeier lattice with root system $A_5^4D_4$ and glue index $72$ (ES result R10) & I & ES's glue code verified; the two codes not compared\\
Jacobian map & Alpöge's map $F:\C^3\to\C^3$ with $\det DF=-2$ and the three-point fibre over $(-\tfrac14,0,0)$ (archive §§4, 7). The same map is in the Yang--Mills workbench and the zeta programme. ES's four-dimensional extension $P$ ($\det DP=-2$) crosses to the zeta programme's conductor & P & $\det DF$, $\det DP$ and the fibre verified (zeta paper \cite{ZetaReader}, Lemma~6.6; \note{21\_}); the conductor maps located\\
ES $\leftrightarrow$ zeta & The split-zero semiring $G(R)=R\sqcup\{\tau\}$ with its Boolean support character (archive §21; ES R01) & I & audited (\note{21\_} §6)\\
ES $\leftrightarrow$ zeta & A finite Weyl phase-space basis & I & ES side verified on $C_2\times C_3$\\
zeta $\to$ ES & NG20--NG21, used by item 20; the SplitZero packages of 13--16 September & D & located\\
ES $\to$ zeta & The 488-page reader's source, vendored into the zeta repository & D & located\\
ES $\to$ zeta & The 124-page ES/Fable reader: solutions as binary quartics, then the Fable cover and a weighted conductor (§\ref{ssec:crosswalk}) & P & the quartic relation proved here (Proposition~\ref{prop:quartic}); the rest located\\
ES $\leftrightarrow$ Collatz & $3(4n+1)+1=4(3n+1)$, i.e.\ $\theta T=4\theta$ (ES R05, diagrams by u/CivQ17) & I & audited\\
Collatz $\leftrightarrow$ ES $\leftrightarrow$ zeta & The homogenisation $(x,h)\mapsto(ax+bh,h)$ (ES X4) & I & ES's diagram read\\
Busy Beaver (new) & Busy-Beaver extrema quoted by the archive ($S(4)=107$ and $\mathrm{Space}(5)=12289$; for the five-state values it cites the formal enumeration of the bbchallenge collaboration) select the shells of archive §§3 and 20; §25 maps extrema to shell carriers. The zeta programme's detectability argument \cite{ZetaReader} uses a different machine for a different purpose & A & located; the two use different machines for different purposes\\
Connes (new) & Supplement §6 reads Connes' primitive intersections; archive §24 opens with Tomita--Takesaki transport. The zeta programme uses Connes' summation map and Meyer's quotient & A & located; no common statement was found\\
Fable (new) & \emph{Fable} names three different objects across the repositories: the three-dimensional Keller map, the $S^6$ construction in older files, and a conductor cover (\note{47a\_} §10). Archive §5 withdraws one Fable construction & --- & disambiguate at each use\\
ES $\leftrightarrow$ YM & A shared term, \emph{gap} (ES X7), and a proposed success projector & A & located; no typed map is written\\
\bottomrule
\end{longtable}
\end{center}

The crossings with typed content are the torus operator, the monodromy matrices, the Keller maps, the split-zero semiring and two identities; each is audited or recorded as located. The other connections are proposals awaiting typed maps, and the archive describes its late Fable, operator, Connes and moonshine material in the same way (p.~4): as a candidate interface, pending an explicit typed map and compatibility check.
